\documentclass[fullscreen=true, bookmarks=false]{beamer}
\usepackage[T2A]{fontenc}
\usepackage[utf8x]{inputenc}
\usepackage[english,russian]{babel}
\usepackage{xcolor}
\usepackage{graphicx}
\usepackage{graphicx}

\title{Санкт-Петербургский Политехнический университет}
\author{Карпунин А.Е.}
\institute{Санкт-Петербургский Политехнический университет}
\date{01.06.13}

\begin{document}
\graphicspath{{noiseimages/}}

\begin{frame}
\transdissolve[duration=0.2]
\titlepage
\end{frame}

\section*{План презентации}


\begin{frame}
\transdissolve[duration=0.1]
\frametitle{Содержание}
\tableofcontents[pausesections]
\end{frame}

\section{Санкт-Петербургский Политехнический университет}
\subsection{Общая информация}
\subsection{Институты}

\begin{frame}
\frametitle{Основная информация}
\textbf{Санкт-Петербургский государственный политехнический университет} (ФГБОУ ВПО «СПбГПУ», полное название — Федеральное государственное бюджетное образовательное учреждение высшего профессионального образования «Санкт-Петербургский государственный политехнический университет», неофициальное название — Политех) — один из крупнейших технических университетов в России. Входит в пятерку лучших технических вузов страны. Входит в группу национальных исследовательских университетов России. Университет включает в себя 12 базовых институтов,факультеты дополнительного образования. Обеспечивает выпуск инженеров, экономистов, менеджеров по 101 специальности, бакалавров и магистров по 34 направлениям науки и техники, аспирантов по 90 научным специальностям.\\

\end{frame}

\begin{frame}
\frametitle{Основная информация}
На 1 января 2007 года в университете по очной форме обучения обучались 18050 человек, всего — более 28000 человек. Профессорско-преподавательский состав включает в себя более 20 академиков и членов-корреспондентов РАН, свыше 500 профессоров, докторов наук.
На базе факультетов СПбГПУ были образованы СПбГУЭФ, СПбГМТУ и СПбВИТУ.
По имени университета названы близлежащая Политехническая улица и станция метро «Политехническая».
\end{frame}

\begin{frame}
\begin{figure}[b]
\noindent\centering{
  \includegraphics[width=0.5\linewidth]{1.png}
  }
  \caption{Герб СПбГПУ}

\end{figure}
\end{frame}

\begin{frame}
\frametitle{Структура}
\textbf{Базовые институты СПбГПУ}: Гуманитарные и экономические, Технические,Физические.\par
\textbf{Факультеты переподготовки специалистов и дополнительного образования}: Факультет переподготовки специалистов (ФПС), Межотраслевой институт переподготовки специалистов (МИПК), Факультет повышения квалификации преподавателей (ФПКП), Переподготовка преподавателей начального и среднего профессионального образования, Повышение Квалификации по программе «Защита государственной тайны»,Курсы повышения квалификации и переподготовки по направлению «Строительство».\par
\textbf{Вечернее отделение}:Информатика, радио и телекоммуникации, электротехника (специальности электромеханического, радиофизического факультетов и факультета технической кибернетики), Экономика, машиностроение и строительство (специальности энергомашиностроительного, механико-машиностроительного факультетов и факультета экономики и менеджмента)
\pause\\
\end{frame}

\end{document} 